\chapter{Forme bilinéaire électrogravitationnelle et équations de réponse croisée}
\label{chap:forma-electrogravitatoria}

\section{Coordonnées internes de source}

Soient \(Z\) la section d'impédance du vide et \((Q,\ell)\) la paire
du chapitre précédent. À une charge \(q\) et une énergie \(E\), nous associons

\[
\xi_q(q)=\sqrt{\frac{Z}{4\pi\hbar}}\,q,
\qquad
\xi_E(E)=\frac{E}{Q}.
\]

Le vecteur interne de source est

\[
X(q,E)=
\begin{pmatrix}\xi_q(q)\\ \xi_E(E)\end{pmatrix},
\qquad
\eta_{1,1}=
\begin{pmatrix}1&0\\0&-1\end{pmatrix}.
\]

La forme \(X_1^{\mathsf T}\eta_{1,1}X_2\) sépare l'orientation électrique et la
clôture gravitationnelle dans un même plan adimensionnel.

\section{Unification statique des potentiels}

\begin{theorem}[Forme électrogravitationnelle]
\label{thm:rm-emg-potential}
Si \(Zc=\varepsilon^{-1}\), \(E_i=m_ic^2\) et
\(Q\ell=\hbar c\), alors
\begin{equation}
\boxed{
V_{12}(r)=\frac{Q\ell}{r}
X(q_1,E_1)^{\mathsf T}\eta_{1,1}X(q_2,E_2)}
\label{eq:rm-emg-potential}
\end{equation}
est exactement la somme du potentiel de Coulomb et du potentiel de Newton.
\end{theorem}

\begin{proof}
Le terme électrique est
\[
\frac{Q\ell}{r}\frac{Zq_1q_2}{4\pi\hbar}
=\frac{Zcq_1q_2}{4\pi r}
=\frac{q_1q_2}{4\pi\varepsilon r}.
\]
Le terme gravitationnel est
\[
-\frac{Q\ell}{r}\frac{E_1E_2}{Q^2}
=-\frac{\ell}{Q}\frac{E_1E_2}{r}
=-\frac{G}{c^4}\frac{m_1m_2c^4}{r}
=-\frac{Gm_1m_2}{r}.
\]
\end{proof}

La force radiale associée prend la forme

\begin{equation}
F_{12}(r)=\frac{Q}{\ell}
\left(\frac{\ell}{r}\right)^2
X_1^{\mathsf T}\eta_{1,1}X_2.
\label{eq:rm-emg-force}
\end{equation}

Le facteur \(Q/\ell=\mathscr T_P\) est la tension universelle ; la forme
adimensionnelle décide de l'intensité et du signe de sa publication.

\section{Cône d'extrémalité des réponses constitutives et classification de ses rayons}

Pour une source individuelle, le discriminant

\[
\Delta_{\rm EMG}(q,E)=E^2-\mathcal Q_e(q)^2,
\qquad
\mathcal Q_e(q)=Q\sqrt{\frac{Z}{4\pi\hbar}}\,q,
\]

induit les trois secteurs

\[
\Delta_{\rm EMG}>0,
\qquad
\Delta_{\rm EMG}=0,
\qquad
\Delta_{\rm EMG}<0.
\]

La condition nulle équivaut à

\begin{equation}
E^2=\mathcal Q_e(q)^2
\quad\Longleftrightarrow\quad
Gm^2=\frac{q^2}{4\pi\varepsilon}.
\label{eq:rm-rn-extremality}
\end{equation}

C'est la condition d'extrémalité électrogravitationnelle dans la normalisation de
Reissner--Nordström. Le trit acquiert ici une réalisation statique :
sous-extrémalité, seuil extrême et surextrémalité sont les trois signes d'une
même forme quadratique.

Pour la charge élémentaire construite par

\[
e^2=\frac{4\pi\alpha_{\rm HMT}\hbar}{Z},
\]

la coordonnée satisfait

\begin{equation}
\boxed{\xi_q(e)^2=\alpha_{\rm HMT}.}
\label{eq:rm-electron-electric-coordinate}
\end{equation}

Cette égalité situe la charge centrale dans le plan des sources sans utiliser la
masse électronique ni importer sa théorie complète.

\section{Coefficient d'Einstein--Maxwell}

Écrivons

\[
\kappa=\frac{8\pi G}{c^4}=8\pi\mathscr C_P.
\]

Dans une convention où le tenseur électromagnétique s'exprime comme
\(T_{\mu\nu}^{\rm EM}=\mu^{-1}\mathcal Q_{\mu\nu}(F,g)\), le coefficient
gravitationnel de la forme quadratique est \(\kappa/\mu\). En utilisant

\[
\ell^2=\frac{G\hbar}{c^3},
\qquad
Z=\mu c,
\qquad
e^2=\frac{4\pi\alpha\hbar}{Z},
\]

on obtient

\begin{equation}
\boxed{
\frac{\hbar^2}{8\pi\ell^2}
\frac{(\kappa/\mu)}{e^2}
=\frac{1}{4\pi\alpha_{\rm HMT}}.}
\label{eq:rm-einstein-maxwell-alpha}
\end{equation}

\begin{proof}
D'abord,
\[
\frac{\kappa}{\mu}
=\frac{8\pi G}{c^4\mu}
=\frac{8\pi\ell^2}{\hbar\mu c}
=\frac{8\pi\ell^2}{\hbar Z}.
\]
En divisant par \(e^2=4\pi\alpha\hbar/Z\) et en multipliant par
\(\hbar^2/(8\pi\ell^2)\), il reste \(1/(4\pi\alpha)\).
\end{proof}

Dans l'interface électrofaible rationalisée,

\[
g^2=\frac{4\pi\alpha}{1-D_A},
\qquad
g'^2=\frac{4\pi\alpha}{D_A},
\]

et donc

\begin{equation}
\boxed{
\frac{\hbar^2}{8\pi\ell^2}
\frac{(\kappa/\mu)}{e^2}
=\frac1{g^2}+\frac1{g'^2}.}
\label{eq:rm-cross-einstein-ew}
\end{equation}

L'égalité impose la commutation des lecteurs de gravité, de vide, de charge et de
mélange de jauge. Dans le graphe actuel, \(e,g,g'\) descendent de
\(\alpha_{\rm HMT}\) ; l'équation est un contrôle croisé et non une deuxième
dérivation indépendante de \(\alpha\).

\section{Polygone des publications de la structure fine}

Le macro-objet dodécaphasique se publie au moyen du système

\begin{equation}
\boxed{
\alpha_{\rm HMT}
=\frac{Ze^2}{4\pi\hbar}
=\frac{ZG_0}{4}
=\frac{Z}{2R_K}
=-\frac{9}{100\pi}\log D_A
=\frac{(1-D_A)g^2}{4\pi}
=\frac{D_Ag'^2}{4\pi}.}
\label{eq:rm-alpha-polygon}
\end{equation}

La première construction causale de \(\alpha\) reste la roue signée,
le sceau et la retenue. Les égalités ultérieures sont des publications
métrologiques ou des contrôles de compatibilité. L'élimination des variables
produit

\[
(1-D_A)g^2=D_Ag'^2
=\frac{Ze^2}{\hbar}
=\pi G_0Z
=\frac{2\pi Z}{R_K}.
\]

La forme duale est

\[
\frac1{g^2}+\frac1{g'^2}
=\frac{\hbar}{Ze^2}
=\frac{R_K}{2\pi Z}
=\frac1{\pi G_0Z}.
\]

Chaque cycle métrologique doit être marqué comme \emph{round trip} : si \(e\),
\(R_K\) ou \(G_0\) ont été définis à partir de \(\alpha\), leur substitution inverse
vérifie la cohérence, mais n'ajoute pas d'information sélective.

\section{Covariance de l'inversion d'action}

Sous la transformation bidirectionnelle,

\[
\frac{e_{\rm pre}^2}{e_{\rm ret}^2}=R_{\rm act},
\qquad
\frac{G_{\rm pre}}{G_{\rm ret}}=R_{\rm act}^{-1}.
\]

Par conséquent,

\begin{equation}
\boxed{G_{\sigma}e_{\sigma}^2=\text{constante de feuillet}.}
\label{eq:rm-ge2-invariant}
\end{equation}

Sous forme constitutive,

\[
\frac{Ze_{\sigma}^2G_{\sigma}}
{c^5t_0^2}
=4\pi\theta_{108}\alpha_{\rm HMT},
\qquad
\theta_{108}=\left(\frac{54}{\pi}\right)^2.
\]

La charge conserve l'orientation de l'action et la réponse gravitationnelle
la compense ; leur produit appartient à la géométrie partagée.

\Needspace{12\baselineskip}
\section{Représentation tensorielle du porteur gravitationnel et réduction transversale de ses degrés de liberté}

La forme bilinéaire \eqref{eq:rm-emg-potential} clôt le secteur statique. La
promotion à une dynamique complète exige un réalisateur

\[
\mathcal R_{\rm EM}:
\mathcal X_{\TPK}^{\rm enr}
\longrightarrow
(e^I{}_{\mu},\omega^{IJ}{}_{\mu},A_\mu)
\]

qui préserve l'identité de Bianchi, les équations de Maxwell, la
conservation du courant et les termes de bord. L'identité
\eqref{eq:rm-cross-einstein-ew} fixe le coefficient que ce réalisateur doit
publier ; elle ne remplace pas ses équations différentielles.

\begin{falsifierbox}
Réfutent la forme électrogravitationnelle : un coefficient différent dans Coulomb ou
Newton en développant \eqref{eq:rm-emg-potential} ; un cône nul qui ne reproduise pas
\eqref{eq:rm-rn-extremality} ; la perte de
\eqref{eq:rm-einstein-maxwell-alpha} ; ou l'utilisation d'une publication
descendante de \(\alpha\) comme sélecteur rétroactif de sa branche dodécaphasique.
\end{falsifierbox}

\paragraph{Structure causale de la confluence électrogravitationnelle.}
La composition démontrée dans ce chapitre s'organise en cinq étapes
typées. L'état enrichi fournit d'abord la phase
\(\alpha_{\rm HMT}\) et son orientation. L'opérateur constitutif publie ensuite
\(e\), \(g\) et \(g'\) dans leurs codomaines respectifs. La clôture gravitationnelle
détermine \(G\) et la chronologie \(\theta_{108}\). Les identités
\eqref{eq:rm-emg-potential}--\eqref{eq:rm-cross-einstein-ew} coordonnent ces
publications au moyen d'un coefficient commun. Enfin,
\(\mathcal R_{\rm EM}\) transporte le résultat vers les champs
\((e^I{}_{\mu},\omega^{IJ}{}_{\mu},A_\mu)\), en préservant les identités
différentielles et les termes de frontière. Cette factorisation fixe
l'antériorité de la construction HMT par rapport à sa réalisation tensorielle et
prépare les relations constitutives développées dans le chapitre suivant.
